\chapter{The Probabilistic Method}

\begin{center}\itshape
If a random colouring has, on average, fewer than one monochromatic clique,\\
then some colouring has none.
\end{center}

An expectation proves existence when it averages a count of failures. The phrase \emph{with high probability} means that the probability of the event tends to $1$ as $n\to\infty$. It is used only after this sentence.

\section{The first moment and the local lemma}

Colour the edges of $K_n$ red or blue, independently and fairly. Let $X$ be the number of monochromatic copies of $K_k$.

\begin{theorem}[Erd\H os]\label{thm:erdos}
If $\binom{n}{k}2^{1-\binom{k}{2}}<1$, some $2$-colouring of $K_n$ has no monochromatic $K_k$. In particular $R(k,k)>2^{k/2}$ for $k\ge3$.
\end{theorem}
\begin{proof}
Linearity gives $\mathbb E[X]=\binom{n}{k}2^{1-\binom{k}{2}}$, since each clique is monochromatic with probability $2^{1-\binom{k}{2}}$. A mean below $1$ forces $P(X=0)>0$. For $n=\lfloor 2^{k/2}\rfloor$ one has $\binom{n}{k}2^{1-\binom{k}{2}}\le 2^{1+k/2}/k!<1$ when $k\ge3$.
\end{proof}

\begin{theorem}[symmetric local lemma]\label{thm:lll}
Let $A_1,\ldots,A_m$ be events with $P(A_i)\le p$. Suppose each event has a set $D_i$ of at most $d$ other events such that $A_i$ is independent of the \emph{joint collection} of all events outside $D_i\cup\{i\}$ (a dependency graph). If $ep(d+1)\le1$, then $P(\bigcap_{i=1}^m A_i^c)>0$.
\end{theorem}
\begin{proof}
If $d=0$, the events are mutually independent, and $p\le1/e<1$ gives $P(\bigcap_iA_i^c)=\prod_i(1-P(A_i))>0$. Assume $d\ge1$ and put $x=1/(d+1)<1$. Since $(1-1/(d+1))^d\ge e^{-1}$ (use $\ln(1+u)\le u$ with $u=1/d$), the hypothesis gives
\[
 p\le\frac1{e(d+1)}\le x(1-x)^d.
\]

We prove by induction on the number of conditioning events that, for every $i\notin S$ and every nonempty conditioning event,
\[
 P\bigl(A_i\mid\bigcap_{j\in S}A_j^c\bigr)\le x.
\]
For $S=\varnothing$, this follows from $p\le x$. For general $S$, split $S=T\cup R$, where $T=S\cap D_i$ has at most $d$ members and $R=S\setminus D_i$. By the dependency-graph condition, $A_i$ is independent of the joint avoidance event for all $j\in R$, so
\[
P(A_i\mid\bigcap_{j\in R}A_j^c)=P(A_i)\le p.
\]
Conditioning further on avoiding $T$ can increase this probability by at most the reciprocal of
\[
P\bigl(\bigcap_{j\in T}A_j^c\mid\bigcap_{j\in R}A_j^c\bigr)
=\prod_{j\in T}P\bigl(A_j^c\mid\text{avoid earlier members of }T\text{ and all }R\bigr)
\ge(1-x)^{|T|}\ge(1-x)^d.
\]
The inductive claim applies in each factor because its conditioning set has at most $|S|-1$ elements. Therefore $P(A_i\mid\text{avoid }S)\le p/(1-x)^d\le x$. The same factors are strictly positive, so all conditioning events used are nonempty. Finally,
\[
P\bigl(\bigcap_{i=1}^m A_i^c\bigr)
=\prod_{i=1}^m P(A_i^c\mid A_1^c,\ldots,A_{i-1}^c)
\ge(1-x)^m>0.
\]
The factor $(1-x)^d$, not the weaker union-bound denominator $1-dx$, is what makes the numerical hypothesis work.
\end{proof}

\section{The second moment}

A mean below $1$ forces a positive chance of zero. A mean that stays at a constant does not. The variable could be zero most of the time and enormous on a rare event, and the average would still be that constant. The second moment forbids that rare spike.

\begin{proposition}[second moment]\label{prop:second}
Let $X$ be a random variable on a finite sample space, with $\mathbb E[X^2]>0$. Then
\[
 P(X\ne0)\ge\frac{(\mathbb E[X])^2}{\mathbb E[X^2]}.
\]
If $X\ge0$ and $\mathbb E[X]>0$, Chebyshev's inequality gives the same information in another shape: $P(X=0)\le\operatorname{Var}(X)/(\mathbb E[X])^2$.
\end{proposition}
\begin{proof}
List the outcomes $\omega$ and their probabilities $p(\omega)$. The Cauchy--Schwarz inequality of Volume~1, Chapter~9, applied to the lists $X(\omega)\sqrt{p(\omega)}$ and $\sqrt{p(\omega)}\,\mathbf 1_{\{X(\omega)\ne0\}}$, gives
\[
 \bigl(\mathbb E[X]\bigr)^2
 =\Bigl(\sum_\omega X(\omega)\sqrt{p(\omega)}\cdot\sqrt{p(\omega)}\,\mathbf 1_{\{X\ne0\}}\Bigr)^2
 \le\Bigl(\sum_\omega X(\omega)^2 p(\omega)\Bigr)\Bigl(\sum_{X(\omega)\ne0}p(\omega)\Bigr).
\]
The second factor is $P(X\ne0)$. For the Chebyshev form, $X\ge0$ and $\mathbb E[X]>0$ turn $\{X=0\}$ into the event $|X-\mathbb E[X]|\ge\mathbb E[X]$.
\end{proof}

The two lower bounds agree after a little algebra. Write $v=\operatorname{Var}(X)/(\mathbb E[X])^2$. Then $\mathbb E[X^2]=(\mathbb E[X])^2(1+v)$, so the ratio in the display is $1/(1+v)$. Chebyshev says only that the probability of being positive is at least $1-v$, which is weaker, and which may be negative. The ratio is the bound to keep.

Three fair coins give a check that does not depend on a graph. Let $X$ be the number of heads. Then $\mathbb E[X]=3/2$ and $\mathbb E[X^2]=3$, because the values $0,1,4,9$ occur with probabilities $1/8,3/8,3/8,1/8$ and $3/8+12/8+9/8=3$. The ratio is $(9/4)/3=3/4$. The true probability $P(X\ne0)=7/8$ is larger, as the inequality requires. Chapter~14 uses the same ratio on triangles.

\section{Euler's formula}

\begin{proposition}[Euler]\label{prop:euler}
A connected plane graph satisfies $v-e+f=2$. A simple connected plane graph with $v\ge3$ has $e\le 3v-6$.
\end{proposition}
\begin{proof}
A tree has $e=v-1$ and one face. Each added edge splits a face and raises $e$ and $f$ by one. For the inequality, $2e\ge 3f$ because every face has at least three edges and every edge meets at most two faces. Substitute $f=e-v+2$.
\end{proof}

A drawing with no crossings is a plane graph. The crossing number in Exercise~\ref{ex:crossing} counts how many crossings a drawing is forced to have. Chapter~14 asks the same question for a random graph: which patterns appear on their own.

\section{Exercises}

\begin{exercise}[$\star$]\label{ex:carowei}
Show that every graph has an independent set of size at least $\sum_v 1/(\deg v+1)$.
\end{exercise}

\begin{exercise}[$\star$]\label{ex:tournament}
Show that some tournament on $n$ vertices has at least $\binom{n}{3}/4$ directed $3$-cycles.
\end{exercise}

\begin{exercise}[$\star\star$]\label{ex:sumfree}
Show that every set of $n$ nonzero integers has a sum-free subset of size greater than $n/3$.
\end{exercise}

\begin{exercise}[$\star\star$]\label{ex:crossing}
Let $\operatorname{cr}(G)$ be the minimum number of crossings in a drawing of a simple graph $G$ with $n\ge3$ vertices and $m$ edges. Using Proposition~\ref{prop:euler}, prove $\operatorname{cr}(G)\ge m-3n$. If $m\ge4n$, prove $\operatorname{cr}(G)\ge m^3/(64n^2)$ by retaining each vertex independently with probability $p=4n/m$.
\end{exercise}

\begin{exercise}[$\star$]\label{ex:propertyb}
Show that a family of fewer than $2^{k-1}$ sets of size $k$ has a $2$-colouring of the ground set with no monochromatic member.
\end{exercise}


\begin{exercise}[$\star\star$]\label{ex:localcolour}
Let $\mathcal F$ be a family of subsets of a finite ground set, each of size $k\ge2$. Assume every set of $\mathcal F$ intersects at most $d$ other sets. Colour the ground-set elements independently red or blue. Prove that $\mathcal F$ has a colouring with no monochromatic member whenever $e(d+1)2^{1-k}\le1$. Explain why the dependency condition is joint independence, not merely pairwise independence of the bad events.
\end{exercise}

\section{Hints}

\begin{hint}{ex:carowei}
In a random order, keep each vertex that precedes its neighbours.
\end{hint}
\begin{hint}{ex:tournament}
A random orientation makes a triple cyclic with probability $1/4$.
\end{hint}
\begin{hint}{ex:sumfree}
Take a prime $p=3q+2$ larger than the set, and a random multiplier. The middle third of $\Z/p\Z$ is sum-free.
\end{hint}
\begin{hint}{ex:crossing}
One deleted edge per crossing leaves a plane graph. A crossing survives in a random subset only if all four endpoints are kept.
\end{hint}
\begin{hint}{ex:propertyb}
A fair colouring makes a $k$-set monochromatic with probability $2^{1-k}$.
\end{hint}
\begin{hint}{ex:localcolour}
A fixed $k$-set is monochromatic with probability $2^{1-k}$. Events attached to disjoint sets are functions of disjoint collections of fair bits.
\end{hint}

\section{Solutions}

\begin{solution}{ex:carowei}
\writeup In a uniform random order, the probability that $v$ is the earliest vertex among itself and its neighbours is $1/(\deg v+1)$. The expected number of such vertices is the sum in the statement, and the set of them is independent.
\end{solution}

\begin{solution}{ex:tournament}
\writeup Orient each edge fairly. Of the eight orientations of a triple, two are cyclic. The expected number of cyclic triples is $\binom{n}{3}/4$.
\end{solution}

\begin{solution}{ex:sumfree}
\writeup Choose a prime $p=3q+2$ larger than every absolute value in the finite set $A$. Such primes exist arbitrarily far out, by an elementary Euclidean argument: if $p_1,\ldots,p_r$ were all primes congruent to $2$ modulo $3$, the integer $M=3p_1\cdots p_r-1$ would be $2$ modulo $3$, divisible by none of those primes, and some prime factor of $M$ would be $2$ modulo $3$ (a product of primes all congruent to $1$ modulo $3$ is $1$ modulo $3$). This contradiction proves the necessary supply of primes. In particular $p$ divides no nonzero element of $A$. For uniform $t\in\{1,\ldots,p-1\}$, keep those $a$ with $ta\bmod p$ in $\{q+1,\ldots,2q+1\}$. That interval is sum-free modulo $p$, because a sum of two of its members lies between $2q+2$ and $4q+2=p+q$, hence outside the interval after reduction. Each fixed $a\ne0$ is kept with probability $(q+1)/(p-1)>1/3$. The expected kept set is larger than $n/3$.
\end{solution}

\begin{solution}{ex:crossing}
\writeup Delete one edge from each crossing. The remainder is plane, so $m-\operatorname{cr}(G)\le 3n-6$ by Proposition~\ref{prop:euler}, and $\operatorname{cr}(G)\ge m-3n$.

Retain each vertex with probability $p=4n/m\le1$. For \emph{every} resulting induced graph, including graphs with fewer than three vertices or several components, the universal inequality $\operatorname{cr}(G')\ge m'-3n'$ holds: for $n'\ge3$ remove crossing edges and apply the planar edge bound, and for $n'\le2$ its right-hand side is nonpositive. Taking expectations requires no ignored error term. A crossing of a fixed original drawing survives only if its four distinct endpoints are retained, with probability $p^4$, while an edge survives with probability $p^2$. Choose a minimum-crossing drawing in general position, with no crossing between adjacent edges away from their common endpoint: such an adjacent crossing can be eliminated by swapping the short crossing tails and smoothing, without introducing new intersections. Also perturb any multiple-edge crossing to distinct pairwise crossings. Every remaining crossing involves exactly four distinct endpoint vertices, and survives vertex deletion with probability $p^4$. Hence the expected number of surviving crossings of this fixed drawing is $p^4\operatorname{cr}(G)$, which is at least the expected crossing number of the induced graph. Consequently $p^4\operatorname{cr}(G)\ge p^2 m-3pn$, and
\[
 \operatorname{cr}(G)\ge\frac{m}{p^2}-\frac{3n}{p^3}=\frac{m^3}{16n^2}-\frac{3m^3}{64n^2}=\frac{m^3}{64n^2}.
\]
\end{solution}

\begin{solution}{ex:propertyb}
\writeup A fair random colouring makes a fixed $k$-set monochromatic with probability $2^{1-k}$. Fewer than $2^{k-1}$ sets give an expectation below $1$.
\end{solution}

\begin{solution}{ex:localcolour}
\idea The only events a monochromatic-set event can depend on are those sharing some colour variable with its set.
\writeup For $F\in\mathcal F$, let $A_F$ mean that all elements of $F$ have the same colour. Both all-red and all-blue have probability $2^{-k}$, so $P(A_F)=2^{1-k}$. If $G$ is disjoint from $F$, $A_G$ depends only on colour variables not used in $A_F$. More strongly, $A_F$ is independent of the \emph{joint family} of all events from sets disjoint from $F$, since all those events together are functions only of colour variables outside $F$. The intersection graph of the sets is therefore a valid dependency graph, with degree at most $d$. Theorem~\ref{thm:lll} applies under $e(d+1)2^{1-k}\le1$ and gives a positive probability that no $A_F$ occurs. Some actual colouring thus has no monochromatic set.
\end{solution}

\section*{Sources and further reading}

\begin{enumerate}
\item P.~Erd\H os, Some remarks on the theory of graphs, \emph{Bull.\ Amer.\ Math.\ Soc.} \textbf{53} (1947), 292--294. Theorem~\ref{thm:erdos}.
\item P.~Erd\H os and L.~Lov\'asz, Problems and results on $3$-chromatic hypergraphs and some related questions, in \emph{Infinite and Finite Sets}, North-Holland, 1975. Theorem~\ref{thm:lll}.
\end{enumerate}

